\section{Related Work}
\label{sec:related_work}

\subsection{IBD Detection Methods}

Identity-by-descent detection from genotype data is a well-established field. GERMLINE~\cite{gusev2009germline} pioneered fast seed-and-extend IBD detection using hashed haplotype segments. IBDseq~\cite{browning2013ibdseq} introduced a likelihood-ratio framework comparing IBD and non-IBD models per site. hap-ibd~\cite{browning2020hap_ibd} combined positional Burrows-Wheeler transforms with HMM-based scoring, achieving biobank-scale performance. iLASH~\cite{shemirani2021ilash} uses locality-sensitive hashing for mega-scale IBD detection in diverse populations, while RaPID~\cite{naseri2019rapid} employs random projections on PBWT for ultra-fast IBD calling. Most recently, PBML~\cite{pbml2026} scales PBWT-based IBD detection to biobank cohorts via kL-SMEMs (set-maximal exact matches of length $\geq k$ shared by $\geq L$ haplotypes), achieving 4.6$\times$ speedup over $\mu$-PBWT. In the pathogen genomics domain, hmmibd-rs~\cite{hmmibdrs2026} reimplements hmmIBD in Rust with Rayon parallelism for \textit{Plasmodium} IBD detection, achieving 100$\times$ speedup over the C original. All these methods operate on phased VCF genotype data, requiring variant calling and statistical phasing as preprocessing steps. No assembly-based or pangenome-based IBD detection tool has been published.

\subsection{Local Ancestry Inference}

Local ancestry inference (LAI) has seen rapid development, with all published methods requiring VCF input. ChromoPainter~\cite{lawson2012chromopainter} pioneered HMM-based chromosome painting using the Li \& Stephens copying model, and SparsePainter~\cite{wang2024sparsepainter} extended this to biobank scale using PBWT-based sparse haplotype matching. RFMix~\cite{maples2013rfmix} uses conditional random forests on local haplotype windows and remains the most widely used reference method. FLARE~\cite{browning2023flare}, from the same group that developed hap-ibd, extends the Li \& Stephens model~\cite{li2003stephens} with composite reference haplotypes for fast and accurate LAI. Gnomix~\cite{hilmarsson2021gnomix} combines a base CRF with a neural network smoother. More recently, Recomb-Mix~\cite{zhou2025recombmix} achieves very fast inference via graph collapsing of recombination graphs, reducing complexity from $O(nm)$ to $O(np)$ where $p$ is the number of populations. Orchestra~\cite{hou2025orchestra} applies CNN and attention layers for 35-population classification but requires days of training on HPC infrastructure.

The only other assembly-based LAI method is ntRoot~\citep{ntroot2025}, which uses k-mer Bloom filters to assign ancestry at $\sim$5\,Mb resolution ($R^2 = 0.96$ against reference panels). However, ntRoot operates at 500$\times$ coarser resolution than \texttt{impop\_k} (5\,Mb vs 10\,kb windows), has no HMM for spatial smoothing, cannot detect IBD, and requires pre-built k-mer databases. Our method, \texttt{impop\_k}, is the first to perform local ancestry inference with per-window HMM resolution directly from pangenome alignments, representing a methodological shift from discrete genotype comparisons to continuous sequence identity observations.

\subsection{Pangenome Alignment Tools}

\begin{sloppypar}
The pangenome infrastructure underlying our approach has matured rapidly.
Minigraph-Cactus~\cite{hickey2024minigraph_cactus} constructs pangenome graphs that capture the full spectrum of structural variation.
The implicit pangenome graph tool \texttt{impg}~\cite{impg} computes pairwise sequence identity via transitive alignment closure using cache-oblivious interval trees.
AGC~\cite{agc} provides efficient compressed storage of assembled genomes.
rustybam~\cite{vollger2026rustybam} offers Rust-based tools for windowed identity computation from PAF/CIGAR data, serving as an independent validation source.
\end{sloppypar}

Recent theoretical work on GBWT haplotype matching~\cite{sanaullah2025gbwt} formalizes set-maximal and long-match queries on pangenome graphs, offering a path toward graph-native IBD detection where long shared paths through the graph directly indicate shared haplotypes.

\subsection{TMRCA-Based Relatedness}

The theoretical convergence on TMRCA-based relatedness (Section~\ref{sec:introduction}) motivates per-window identity as a local eGRM. All prior ``full'' GRM methods~\cite{fan2022egrm,lehmann2026branch_grm,gesi2025} require ARG inference; our approach computes $I = 1 - 2\mu t$ directly from assemblies, bypassing both VCF generation and ARG reconstruction. A transformer for TMRCA prediction~\citep{cxt2026} independently validates the identity$\to$TMRCA mapping. The dimensional reduction from $m$ reference haplotypes to $k$ populations---analogous to the Li \& Stephens model~\cite{li2003stephens} at the population level---requires fewer reference haplotypes for stable estimation while sacrificing multi-site haplotype information.
